\section{Related work}
\label{sec:related_work}
\begin{figure*}[t]
\centering
  \includegraphics[width=0.95\textwidth]{figs/totalpipelinev4.pdf}
  \caption{ \textbf{Proposed DynamicVGGT training framework.} Given a sequence of multi-view images $\{V_1,V_2,V_3\}$, the model first encodes them using a pretrained DINOv2 backbone to extract patch tokens and camera tokens for each view, while motion tokens are initialized as learnable parameters that encode temporal priors. 
  The patch and camera tokens are processed by the Alternating-Attention (AA) blocks to model intra-frame spatial geometry, 
  whereas the Motion-aware Temporal Attention (MTA) blocks operate in parallel to model inter-frame temporal dependencies using the motion tokens.
  The resulting temporal features ${TA}$ are then fed into a Dynamic 3D Gaussian Head (DGSHead) for dynamic 3DGS reconstruction and a Future Point Head for future point prediction.
  }
  \label{fig:pipeline}
\end{figure*}
\subsection{Feed-Forward Visual Geometry Learning.}
% [介绍 DUSt3R, VGGT 等前馈几何模型的发展与局限]
Feed-forward 3D reconstruction models~\cite{wang2024dust3r, wang2025vggt, jiang2025anysplat} aim to recover the geometry of static scenes directly from images under a temporal invariance assumption, providing robust 3D priors for downstream tasks~\citep{chen2024vpl, li2025papl, li2025constrained, liartdeco, li2026ec, gaodeep, gaogood}.
Unlike traditional multi-view geometry pipelines~\cite{mur2015orb, mur2017orb} that relied on optimization-based correspondence matching, these learning-based methods predicted depth, camera pose, or dense 3D point maps in an end-to-end manner. 
Representative frameworks such as DUSt3R~\cite{wang2024dust3r} had demonstrated that transformer-based architectures can effectively learn direct mappings from image pixels to 3D coordinate fields. Subsequent researches had expanded this paradigm to multi-view and sequential settings, integrating camera pose estimation, depth prediction, and correspondence learning into unified architectures. 
Among these, VGGT~\cite{wang2025vggt} further enhanced the feed-forward formulation by introducing alternating attention mechanisms across spatial and temporal dimensions, achieving joint prediction of multiple geometric quantities through a single, shared model. 
Anysplat~\cite{jiang2025anysplat} incorporated 3DGS with Feed-Forward Model to achieve high-fidelity reconstruction. Recent extensions of feed-forward visual geometry models have begun to incorporate temporal modeling for dynamic reconstruction. Approaches such as MoVieS\cite{lin2025movies} and StreamVGGT\cite{zhuo2025streaming} extend static frameworks like VGGT to handle sequential inputs, but they are primarily designed for indoor environments.
Despite these advances, existing 3D feed-forward methods struggled with reconstruct the large scale, dynamic autonomous driving scenes, motivating the need for generalizable 4D feed-forward frameworks capable of capturing scene dynamics over time.


\subsection{3DGS Reconstruction for Driving Scenes.}
% [介绍动态 3DGS 方法在自动驾驶上的相关工作， DynamicVGGT做自动驾驶数据集]
Building photorealistic reconstructions of dynamic urban scenes is crucial for autonomous driving, as it supports closed-loop training and evaluation under realistic motion patterns. Consequently, recent research~\cite{yan2024street} has shifted focus from static to dynamic scene reconstruction. However, most existing methods~\cite{chen2024omnire, zhou2024hugs} rely heavily on dense annotations, which are expensive to obtain and limit scalability. Furthermore, these approaches typically depend on per-scene optimization, making it difficult to exploit large-scale data priors and leading to slow reconstruction speeds.
Recent advances in feed-forward reconstruction, such as STORM~\cite{yang2024storm} and DrivingForward~\cite{tian2025drivingforward}, demonstrate the potential for fast and generalizable 4D scene recovery without per-scene optimization.
STORM introduces a feed-forward pipeline for dynamic scene reconstruction and editing, achieving high-quality results but still relying on calibrated multi-view inputs.
DrivingForward~\cite{tian2025drivingforward} presents a feed-forward 3D Gaussian Splatting framework for driving scenes with flexible surround-view inputs, jointly training pose, depth, and Gaussian heads to infer camera poses and dense geometry without using depth ground truth or provided extrinsics. Our DynamicVGGT generalizes feed-forward reconstruction to real-world autonomous driving scenes, jointly modeling geometry and motion without camera parameters or dense annotations.


